% !PRINTED 305
\weberpage{305}

\begin{center}
\textbf{§ 82. \gesperrt{Teiler von $\mathfrak{L}_0$ vom Index $p$ für $p = 5, 7, 11$.}}
\end{center}

Es bleibt die Möglichkeit übrig, daß für $p = 5, 7, 11$ Teiler von $\mathfrak{L}_0$ vom Index $p$ existieren.

1. $p = 5$. Wir untersuchen zunächst den Fall $p = 5$. Da $m$ und $m - 1$ nicht kongruent 0 sein können, so bleiben für $m$ die Annahmen
\[
m \equiv 2, 3, 4 \quad (\mathrm{mod}\ p).
\]
Ersetzen wir aber die Substitution $W$ § 81, (21) durch
\weq{1}{WB = \begin{pmatrix} -\xi\varrho, & \xi \\ -\xi^{-1}m, & \xi^{-1}\varrho^{-1}(m-1) \end{pmatrix},}
wodurch $m$ in $1 - m$ übergeht, so kommt der Fall $m \equiv 4$ auf den Fall $m \equiv 2$ zurück.

Für $m \equiv 2$ ist aber (23), § 81 nicht erfüllt, und von den Kongruenzen (22) ist keine möglich, da $\pm 2$ Nichtreste von 5 sind und die linken Seiten von (22), 3., 4. sich für $m = 2$ auf $\xi^4 - 1$ reduzieren, was für jedes von Null verschiedene $\xi$ durch 5 teilbar ist. Es bleibt also nur übrig, daß
\weq{2}{m \equiv 3, \quad m' \equiv 3 \quad (\mathrm{mod}\ 5)}
und die beiden Systeme $W, W'$ [§ 81, (16), (17)] können sich nur dadurch voneinander unterscheiden, daß $\alpha^{-1}\beta$ und also auch $\varrho$ in dem einen quadratischer Rest, in dem anderen quadratischer Nichtrest von 5 ist. Die gesuchte Gruppe besteht also, falls sie existiert, aus den 12 in den drei Formen
\weq{3}{\begin{pmatrix} \xi, & 0 \\ 0, & \xi^{-1} \end{pmatrix}, \quad \begin{pmatrix} 0, & \xi \\ -\xi^{-1}, & 0 \end{pmatrix}, \quad \begin{pmatrix} \xi, & \eta \\ 2\eta^{-1}, & 3\xi^{-1} \end{pmatrix}}
enthaltenen Substitutionen, worin
\[
\xi \equiv 1, 2, \quad \eta \equiv \pm 1, \pm 2 \quad (\mathrm{mod}\ 5)
\]
zu setzen ist.

2. $p = 7, 11$.

Da $-1$ für $p = 7, 11$ unter den quadratischen Nichtresten zu finden ist, so gehört die zusammengesetzte Substitution $WB$, (1), zu den $W'$ [weil der quadratische Charakter der in (18), § 81 mit $\varrho$ bezeichneten Größe in $W$ und $WB$ der entgegengesetzte ist]. Demnach ist nach (1) und § 81, (19), (21) $\beta'\gamma' \equiv -m \equiv m' - 1$, also:
\weq{4}{m + m' \equiv 1 \quad (\mathrm{mod}\ p).}
